\section{Evaluation}

% A general evaluation of project and/or product then going into more specifics of
% how we solved each Research question (RQ)

% NOTE! It might be beneficial to merge some of the RQ sections 
% together as they could have been solved by the same part of the project.

As the project comes to an end its important to evaluate how well we succeded in our hypothesis and research questions. The project as a whole
has shown to be a success in terms of delivering a functional PCG system for MiniDungeons 2 that is capable of generating dungeons based 
on player experience. The issue has been getting real player experience, and therefore using a simple heuristics Agent to simulate player experience 
was used instead. 

\subsection{Hypothesis Evaluation}
Its clear that the hypothesis of being able to generate dungeons based on player experience has been met. The PCG system is capable 
of generating dungeons based on the experience of the player, and the dungeons are tailored to the players experience, when it begins a 
new test run. It would have been better to make the dynmaic difficulty adjustment part of the run time instead of only at the start of 
a new run. This would have allowed for a more dynamic experience, and made it more noticable for the player.

\subsection{Evaluation of RQ 1 \& 2}
Both research questions have been answered in the project. The first research question was about the ability to use a genetic algorithm to generate
dungeon layouts that achieve a balance between challenge and skill to create a good flow in a procedural dungeon crawler. This has been
achieved through the implementation of the genetic algorithm, which is capable of generating dungeons based on the player experience.
The second research question was about how procedural generation methods (Genetic Algorithm vs. Reinforcement Learning Algorithms) interact to influence 
player experience, engagement and gameplay outcomes. This has been explored through the implementation of the genetic algorithm, and the comparison with
Proximal Policy Optimization (PPO) reinforcement learning algorithm. The results show that the genetic algorithm, while being slower, is capable of 
generating dungeons that are tailored to the player experience, while the PPO algorithm if computational faster, it had a harder time to 
generate dungeons that created a good flow for the player.